% Options for packages loaded elsewhere
\PassOptionsToPackage{unicode}{hyperref}
\PassOptionsToPackage{hyphens}{url}
\PassOptionsToPackage{dvipsnames,svgnames,x11names}{xcolor}
%
\documentclass[
  11pt,
]{article}

\usepackage{amsmath,amssymb}
\usepackage{iftex}
\ifPDFTeX
  \usepackage[T1]{fontenc}
  \usepackage[utf8]{inputenc}
  \usepackage{textcomp} % provide euro and other symbols
\else % if luatex or xetex
  \usepackage{unicode-math}
  \defaultfontfeatures{Scale=MatchLowercase}
  \defaultfontfeatures[\rmfamily]{Ligatures=TeX,Scale=1}
\fi
\usepackage{lmodern}
\ifPDFTeX\else  
    % xetex/luatex font selection
    \setmainfont[]{Arial}
\fi
% Use upquote if available, for straight quotes in verbatim environments
\IfFileExists{upquote.sty}{\usepackage{upquote}}{}
\IfFileExists{microtype.sty}{% use microtype if available
  \usepackage[]{microtype}
  \UseMicrotypeSet[protrusion]{basicmath} % disable protrusion for tt fonts
}{}
\makeatletter
\@ifundefined{KOMAClassName}{% if non-KOMA class
  \IfFileExists{parskip.sty}{%
    \usepackage{parskip}
  }{% else
    \setlength{\parindent}{0pt}
    \setlength{\parskip}{6pt plus 2pt minus 1pt}}
}{% if KOMA class
  \KOMAoptions{parskip=half}}
\makeatother
\usepackage{xcolor}
\usepackage[a4paper, margin=0.8in]{geometry}
\setlength{\emergencystretch}{3em} % prevent overfull lines
\setcounter{secnumdepth}{5}
% Make \paragraph and \subparagraph free-standing
\makeatletter
\ifx\paragraph\undefined\else
  \let\oldparagraph\paragraph
  \renewcommand{\paragraph}{
    \@ifstar
      \xxxParagraphStar
      \xxxParagraphNoStar
  }
  \newcommand{\xxxParagraphStar}[1]{\oldparagraph*{#1}\mbox{}}
  \newcommand{\xxxParagraphNoStar}[1]{\oldparagraph{#1}\mbox{}}
\fi
\ifx\subparagraph\undefined\else
  \let\oldsubparagraph\subparagraph
  \renewcommand{\subparagraph}{
    \@ifstar
      \xxxSubParagraphStar
      \xxxSubParagraphNoStar
  }
  \newcommand{\xxxSubParagraphStar}[1]{\oldsubparagraph*{#1}\mbox{}}
  \newcommand{\xxxSubParagraphNoStar}[1]{\oldsubparagraph{#1}\mbox{}}
\fi
\makeatother


\providecommand{\tightlist}{%
  \setlength{\itemsep}{0pt}\setlength{\parskip}{0pt}}\usepackage{longtable,booktabs,array}
\usepackage{calc} % for calculating minipage widths
% Correct order of tables after \paragraph or \subparagraph
\usepackage{etoolbox}
\makeatletter
\patchcmd\longtable{\par}{\if@noskipsec\mbox{}\fi\par}{}{}
\makeatother
% Allow footnotes in longtable head/foot
\IfFileExists{footnotehyper.sty}{\usepackage{footnotehyper}}{\usepackage{footnote}}
\makesavenoteenv{longtable}
\usepackage{graphicx}
\makeatletter
\newsavebox\pandoc@box
\newcommand*\pandocbounded[1]{% scales image to fit in text height/width
  \sbox\pandoc@box{#1}%
  \Gscale@div\@tempa{\textheight}{\dimexpr\ht\pandoc@box+\dp\pandoc@box\relax}%
  \Gscale@div\@tempb{\linewidth}{\wd\pandoc@box}%
  \ifdim\@tempb\p@<\@tempa\p@\let\@tempa\@tempb\fi% select the smaller of both
  \ifdim\@tempa\p@<\p@\scalebox{\@tempa}{\usebox\pandoc@box}%
  \else\usebox{\pandoc@box}%
  \fi%
}
% Set default figure placement to htbp
\def\fps@figure{htbp}
\makeatother

\makeatletter
\@ifpackageloaded{caption}{}{\usepackage{caption}}
\AtBeginDocument{%
\ifdefined\contentsname
  \renewcommand*\contentsname{Table of contents}
\else
  \newcommand\contentsname{Table of contents}
\fi
\ifdefined\listfigurename
  \renewcommand*\listfigurename{List of Figures}
\else
  \newcommand\listfigurename{List of Figures}
\fi
\ifdefined\listtablename
  \renewcommand*\listtablename{List of Tables}
\else
  \newcommand\listtablename{List of Tables}
\fi
\ifdefined\figurename
  \renewcommand*\figurename{Figure}
\else
  \newcommand\figurename{Figure}
\fi
\ifdefined\tablename
  \renewcommand*\tablename{Table}
\else
  \newcommand\tablename{Table}
\fi
}
\@ifpackageloaded{float}{}{\usepackage{float}}
\floatstyle{ruled}
\@ifundefined{c@chapter}{\newfloat{codelisting}{h}{lop}}{\newfloat{codelisting}{h}{lop}[chapter]}
\floatname{codelisting}{Listing}
\newcommand*\listoflistings{\listof{codelisting}{List of Listings}}
\makeatother
\makeatletter
\makeatother
\makeatletter
\@ifpackageloaded{caption}{}{\usepackage{caption}}
\@ifpackageloaded{subcaption}{}{\usepackage{subcaption}}
\makeatother

\usepackage{bookmark}

\IfFileExists{xurl.sty}{\usepackage{xurl}}{} % add URL line breaks if available
\urlstyle{same} % disable monospaced font for URLs
\hypersetup{
  pdftitle={03\_MBOA\_accessions},
  pdfauthor={Henry Janse van Rensburg},
  colorlinks=true,
  linkcolor={blue},
  filecolor={Maroon},
  citecolor={Blue},
  urlcolor={Blue},
  pdfcreator={LaTeX via pandoc}}


\title{03\_MBOA\_accessions}
\author{Henry Janse van Rensburg}
\date{2026-11-16}

\begin{document}
\maketitle

\renewcommand*\contentsname{Table of contents}
{
\hypersetup{linkcolor=}
\setcounter{tocdepth}{2}
\tableofcontents
}

\section{Background}\label{background}

We conducted this experiment to study the allelopathic effect of MBOA on
Arabidopsis growth. We looked at both root inhibition on plates and
growth inhibition in soil.

\section{Materials and methods}\label{materials-and-methods}

\begin{itemize}
\item
  \textbf{Plant material}: Arabidopsis thaliana
\item
  \textbf{Genotypes}: Col-0 and the most extreme accessions
\item
  \textbf{Treatment}: MBOA
\item
  \textbf{Phenotypes}: Root length, plant area
\item
  \textbf{Accessions}: See list of plant genotypes used in GWAS
\item
  \textbf{Purpose}: We performed this study to test the sensitvity of
  the most extreme accessions to MBOA.
\end{itemize}

\section{Important results}\label{important-results}

MBOA show similar effects to all the accessions, irrespective of their
growth feedbacks

\section{Setup and Load Libraries}\label{setup-and-load-libraries}

\section{Import Phenotype Data}\label{import-phenotype-data}

\subsubsection{Helper Functions}\label{helper-functions}

\begin{itemize}
\tightlist
\item
  \textbf{Purpose}: We defined two helper functions to display the
  number of replicates below each bar in the boxplot and to create
  colored facets at the bottom of the plot.
\end{itemize}

\section{Anova test of Root Length response to
MBOA}\label{anova-test-of-root-length-response-to-mboa}

\begin{verbatim}
                     Df Sum Sq Mean Sq F value   Pr(>F)    
treatment             1  33.20   33.20 123.622  < 2e-16 ***
accession            12 105.93    8.83  32.870  < 2e-16 ***
treatment:accession  12  10.72    0.89   3.327 9.38e-05 ***
Residuals           941 252.72    0.27                     
---
Signif. codes:  0 '***' 0.001 '**' 0.01 '*' 0.05 '.' 0.1 ' ' 1
\end{verbatim}

\begin{longtable}[]{@{}
  >{\raggedright\arraybackslash}p{(\linewidth - 12\tabcolsep) * \real{0.2206}}
  >{\raggedright\arraybackslash}p{(\linewidth - 12\tabcolsep) * \real{0.1471}}
  >{\raggedleft\arraybackslash}p{(\linewidth - 12\tabcolsep) * \real{0.1471}}
  >{\raggedleft\arraybackslash}p{(\linewidth - 12\tabcolsep) * \real{0.1471}}
  >{\raggedleft\arraybackslash}p{(\linewidth - 12\tabcolsep) * \real{0.0588}}
  >{\raggedleft\arraybackslash}p{(\linewidth - 12\tabcolsep) * \real{0.1324}}
  >{\raggedleft\arraybackslash}p{(\linewidth - 12\tabcolsep) * \real{0.1471}}@{}}
\caption{Pairwise Comparisons of Root Length by Treatment and
Accessions}\tabularnewline
\toprule\noalign{}
\begin{minipage}[b]{\linewidth}\raggedright
contrast
\end{minipage} & \begin{minipage}[b]{\linewidth}\raggedright
accession
\end{minipage} & \begin{minipage}[b]{\linewidth}\raggedleft
estimate
\end{minipage} & \begin{minipage}[b]{\linewidth}\raggedleft
SE
\end{minipage} & \begin{minipage}[b]{\linewidth}\raggedleft
df
\end{minipage} & \begin{minipage}[b]{\linewidth}\raggedleft
t.ratio
\end{minipage} & \begin{minipage}[b]{\linewidth}\raggedleft
p.value
\end{minipage} \\
\midrule\noalign{}
\endfirsthead
\toprule\noalign{}
\begin{minipage}[b]{\linewidth}\raggedright
contrast
\end{minipage} & \begin{minipage}[b]{\linewidth}\raggedright
accession
\end{minipage} & \begin{minipage}[b]{\linewidth}\raggedleft
estimate
\end{minipage} & \begin{minipage}[b]{\linewidth}\raggedleft
SE
\end{minipage} & \begin{minipage}[b]{\linewidth}\raggedleft
df
\end{minipage} & \begin{minipage}[b]{\linewidth}\raggedleft
t.ratio
\end{minipage} & \begin{minipage}[b]{\linewidth}\raggedleft
p.value
\end{minipage} \\
\midrule\noalign{}
\endhead
\bottomrule\noalign{}
\endlastfoot
control - MBOA & Bozen-1.2 & 0.0571634 & 0.1322528 & 941 & 0.432228 &
0.6656747 \\
control - MBOA & BRR60 & 0.9331417 & 0.1247897 & 941 & 7.477715 &
0.0000000 \\
control - MBOA & Col-0 & 0.4045364 & 0.0964189 & 941 & 4.195613 &
0.0000298 \\
control - MBOA & Dra2-1 & 0.2115331 & 0.1289406 & 941 & 1.640547 &
0.1012257 \\
control - MBOA & DraIV 3-7 & 0.2317944 & 0.0964113 & 941 & 2.404225 &
0.0163992 \\
control - MBOA & Giffo-1 & 0.3810455 & 0.1221966 & 941 & 3.118300 &
0.0018743 \\
control - MBOA & Hadd-3 & 0.2958328 & 0.1082902 & 941 & 2.731852 &
0.0064159 \\
control - MBOA & IP-Hue-3 & 0.3460893 & 0.1248947 & 941 & 2.771049 &
0.0056977 \\
control - MBOA & Kastel-1 & 0.3927718 & 0.1214124 & 941 & 3.235021 &
0.0012587 \\
control - MBOA & Lm-2 & 0.4205134 & 0.1709578 & 941 & 2.459749 &
0.0140823 \\
control - MBOA & Rak-2 & 0.2539568 & 0.1240849 & 941 & 2.046638 &
0.0409711 \\
control - MBOA & T840 & 0.7967438 & 0.1438392 & 941 & 5.539126 &
0.0000000 \\
control - MBOA & Tac-0 & 0.3906219 & 0.1266393 & 941 & 3.084523 &
0.0020983 \\
\end{longtable}

\pandocbounded{\includegraphics[keepaspectratio]{03_MBOA_Accessions_files/figure-pdf/Pairwise comparison-1.pdf}}

\begin{itemize}
\tightlist
\item
  \textbf{Result}: MBOA inhibits root growth in all accessions,
  irrespective of their growth feedback. The root length is
  significantly reduced in the MBOA treatment compared to the control.
\end{itemize}

\section{Treatment response}\label{treatment-response}

\pandocbounded{\includegraphics[keepaspectratio]{03_MBOA_Accessions_files/figure-pdf/Response-1.pdf}}

\section{Filter Data for Specific
Genotypes}\label{filter-data-for-specific-genotypes}

\begin{longtable}[]{@{}rlrlrr@{}}
\caption{Extreme contrasting Accessions at DAG = 15}\tabularnewline
\toprule\noalign{}
Snapshot\_ID\_Tag & genotype & genotype\_id & bx & dag &
plant\_region\_area \\
\midrule\noalign{}
\endfirsthead
\toprule\noalign{}
Snapshot\_ID\_Tag & genotype & genotype\_id & bx & dag &
plant\_region\_area \\
\midrule\noalign{}
\endhead
\bottomrule\noalign{}
\endlastfoot
24540 & DraIV 3-7 & 5921 & - & 15 & 3835 \\
24140 & Tac-0 & 7350 & + & 15 & 61168 \\
23912 & T840 & 6136 & + & 15 & 38120 \\
\end{longtable}

\begin{itemize}
\tightlist
\item
  \textbf{Purpose}: We filtered the \texttt{plant\_region\_ares\_GWAS}
  data frame to include only the rows where \texttt{dag} is equal to 15
  and the \texttt{genotype\_id} is one of the selected genotypes.
\end{itemize}

\subsection{Calculate and Merge Average
Values}\label{calculate-and-merge-average-values}

\begin{longtable}[]{@{}rrrr@{}}
\caption{Rosette area of extreme contrasting Accessions at DAG =
15}\tabularnewline
\toprule\noalign{}
genotype\_id & + & - & response \\
\midrule\noalign{}
\endfirsthead
\toprule\noalign{}
genotype\_id & + & - & response \\
\midrule\noalign{}
\endhead
\bottomrule\noalign{}
\endlastfoot
506 & 18687.67 & 53433.833 & -34746.17 \\
5867 & 49635.50 & 15475.500 & 34160.00 \\
5921 & 46544.67 & 9785.667 & 36759.00 \\
\end{longtable}

\begin{longtable}[]{@{}rrrr@{}}
\caption{Root lenght of extreme contrasting Accessions in response to
MBOA}\tabularnewline
\toprule\noalign{}
accession\_id & control & MBOA & treatment\_response \\
\midrule\noalign{}
\endfirsthead
\toprule\noalign{}
accession\_id & control & MBOA & treatment\_response \\
\midrule\noalign{}
\endhead
\bottomrule\noalign{}
\endlastfoot
506 & 2.557561 & 1.624419 & -0.9331417 \\
5867 & 1.422431 & 1.210898 & -0.2115331 \\
5921 & 1.413345 & 1.181551 & -0.2317944 \\
\end{longtable}

\begin{itemize}
\tightlist
\item
  \textbf{Purpose}: We calculated the average plant region area for each
  genotype and BX treatment and then calculated the response by
  subtracting the BX minus values from the BX plus values. We also
  calculated the average root length for each accession and treatment
  and calculated the treatment response by subtracting the MBOA values
  from the control values.
\end{itemize}

\subsection{Merge Data and Calculate
Correlation}\label{merge-data-and-calculate-correlation}

\begin{verbatim}
[1] "Correlation:  0.49"
\end{verbatim}

\begin{verbatim}
[1] "P-value:  0.0929578480866173"
\end{verbatim}

\begin{itemize}
\item
  \textbf{Purpose}: We merged the \texttt{average\_root\_length\_wide}
  and \texttt{average\_area\_wide} data frames on the
  \texttt{accession\_id} column and assigned colors based on the
  response value. We then calculated the correlation between the
  treatment response and the response values and the associated p-value.
\item
  \textbf{Result}: The correlation between the treatment response and
  the response values is 0.485, with a p-value of 0.092. This indicates
  that there is no correlation between the two variables.
\end{itemize}

\subsubsection{Plot Correlation and Treatment
Response}\label{plot-correlation-and-treatment-response}

\begin{figure}[H]

{\centering \pandocbounded{\includegraphics[keepaspectratio]{03_MBOA_Accessions_files/figure-pdf/Fig-MBOA_response-1.pdf}}

}

\caption{Extended Data Fig. 3. Direct allelopathic effect of
benzoxazinoids on the growth of Arabidopsis. (a) Sensitivity to MBOA in
terms of root length response (root length on control -- root length on
MBOA) of Col-0 and the accessions with the strongest positive or
negative microbiome feedback. Accessions are colour-coded as to whether
they showed a positive (blue) or negative (orange) microbiome feedback
on the BX-conditioned soil microbiome in the GWAS experiment. Plants
were grown for 10 days on ½ MS media in the presence of 0.25 mM MBOA or
control (DMSO). The number of seedlings (n) behind each accession and
treatment is given in brackets after the accession name on the y axis as
n MBOA / n control. Boxplot insert displays the root length of Col-0 on
the control or MBOA treatment as an example. Each point is an individual
seedling, i.e.~a biological replicate, distributed over two replicate
plates per treatment. Boxplot shows the median (centre line) and the
25th and 75th percentiles (box bounds); whiskers extend to the most
extreme value within 1.5 × the interquartile range of the nearest box
bound, and all seedlings are plotted as points. Root length was analysed
by two-way ANOVA. MBOA was compared with control within each accession
by two-sided t-tests on the estimated marginal means of this model (one
comparison per accession; unadjusted P values shown on the plot). (b)
Pearson's correlation analysis between MBOA sensitivity in (a) and the
microbiome feedback of the GWAS experiment at 15 days after germination
across the different accessions. The scatter plot displays individual
data points coloured by microbiome feedback, with positive responses
shown in blue and negative responses in orange. Each point is one
accession, its MBOA response being the mean over the seedlings in (a)
and its microbiome feedback the mean over three replicate pots per soil.
A linear regression line (in red) illustrates the trend between the two
variables, surrounded by its pointwise 95\% confidence band (grey
shading). Pearson's correlation coefficient (r) and the exact two-sided
P value are shown on the plot . Dataset 2.}

\end{figure}%

\begin{figure}[H]

{\centering \pandocbounded{\includegraphics[keepaspectratio]{03_MBOA_Accessions_files/figure-pdf/Fig-MBOA_response-2.pdf}}

}

\caption{Extended Data Fig. 3. Direct allelopathic effect of
benzoxazinoids on the growth of Arabidopsis. (a) Sensitivity to MBOA in
terms of root length response (root length on control -- root length on
MBOA) of Col-0 and the accessions with the strongest positive or
negative microbiome feedback. Accessions are colour-coded as to whether
they showed a positive (blue) or negative (orange) microbiome feedback
on the BX-conditioned soil microbiome in the GWAS experiment. Plants
were grown for 10 days on ½ MS media in the presence of 0.25 mM MBOA or
control (DMSO). The number of seedlings (n) behind each accession and
treatment is given in brackets after the accession name on the y axis as
n MBOA / n control. Boxplot insert displays the root length of Col-0 on
the control or MBOA treatment as an example. Each point is an individual
seedling, i.e.~a biological replicate, distributed over two replicate
plates per treatment. Boxplot shows the median (centre line) and the
25th and 75th percentiles (box bounds); whiskers extend to the most
extreme value within 1.5 × the interquartile range of the nearest box
bound, and all seedlings are plotted as points. Root length was analysed
by two-way ANOVA. MBOA was compared with control within each accession
by two-sided t-tests on the estimated marginal means of this model (one
comparison per accession; unadjusted P values shown on the plot). (b)
Pearson's correlation analysis between MBOA sensitivity in (a) and the
microbiome feedback of the GWAS experiment at 15 days after germination
across the different accessions. The scatter plot displays individual
data points coloured by microbiome feedback, with positive responses
shown in blue and negative responses in orange. Each point is one
accession, its MBOA response being the mean over the seedlings in (a)
and its microbiome feedback the mean over three replicate pots per soil.
A linear regression line (in red) illustrates the trend between the two
variables, surrounded by its pointwise 95\% confidence band (grey
shading). Pearson's correlation coefficient (r) and the exact two-sided
P value are shown on the plot . Dataset 2.}

\end{figure}%

\begin{itemize}
\item
  \textbf{Purpose}: We plotted the correlation between the treatment
  response and the response values and the treatment response of each
  accession. The correlation plot shows the relationship between the
  treatment response and the response values, while the treatment
  response plot shows the response of each accession to the treatment.
\item
  \textbf{Purpose}: We combined the two plots into a single plot for
  better visualization and saved the combined plot as a PDF.
\end{itemize}




\end{document}
